% Useful variables
\newcommand{\DocMainTitle}{Documenter LaTeX Showcase}
\newcommand{\DocVersion}{1.2.3}
\newcommand{\DocAuthors}{}
\newcommand{\JuliaVersion}{1.11.1}

% ---- Insert preamble
%% Default preamble BEGIN
\documentclass[oneside]{memoir}

\usepackage{./documenter}
\usepackage{./custom}

%% Title Page
\title{
    {\HUGE \DocMainTitle }\\
    {\Large \DocVersion }
}
\author{ \DocAuthors }

%% TOC settings
% -- TOC depth
%   value: [part, chapter, section, subsection,
%           subsubsection, paragraph, subparagraph]
\settocdepth{section}  % show "part+chapter+section" in TOC
% -- TOC spacing
%   ref: https://tex.stackexchange.com/questions/60317/toc-spacing-in-memoir
%   doc: memoir/memman.pdf
%       - Figure 9.2: Layout of a ToC
%       - Table 9.3: Value of K in macros for styling entries
\makeatletter
% {part} to {chaper}
\setlength{\cftbeforepartskip}{1.5em \@plus \p@}
% {chaper} to {chaper}
\setlength{\cftbeforechapterskip}{0.0em \@plus \p@}
% Chapter num to chapter title spacing (Figure 9.2@memman)
\setlength{\cftchapternumwidth}{2.5em \@plus \p@}
% indent before section number
\setlength{\cftsectionindent}{2.5em \@plus \p@}
% Section num to section title spacing (Figure 9.2@memman)
\setlength{\cftsectionnumwidth}{4.0em \@plus \p@}
\makeatother

%% Main document begin
\begin{document}

\frontmatter
\maketitle
\clearpage
\tableofcontents

\mainmatter
%% preamble END



\part{Showcase}


\chapter{Style showcase for \LaTeX}



\label{5852870456629818970}{}


This page showcases the various page elements that are supported by Documenter.



\section{Table of contents}



\label{11192443102986465861}{}


A table of contents can be generated with an \hyperlinkref{5495928211144802529}{\texttt{@contents} block}. The one for this page renders as


\begin{itemize}
\item \hyperlinkref{5852870456629818970}{Style showcase for \LaTeX}
\begin{itemize}
\item \hyperlinkref{11192443102986465861}{Table of contents}
\item \hyperlinkref{13015211943483912510}{Basic Markdown}
\item \hyperlinkref{14660894509460337656}{Code blocks}
\item \hyperlinkref{9406902154236239312}{Mathematics}
\item \hyperlinkref{5224713065671704740}{Images}
\item \hyperlinkref{1916796393499649241}{Admonitions}
\item \hyperlinkref{932242378642137797}{Lists}
\item \hyperlinkref{6147076236960597988}{Tables}
\item \hyperlinkref{16845496644619434430}{Footnotes}
\item \hyperlinkref{5680510222270408269}{Headings}
\item \hyperlinkref{10597583416623881668}{Doctesting example}
\item \hyperlinkref{3545462789522369084}{Running interactive code}
\item \hyperlinkref{5969831508618549664}{Setup blocks}
\item \hyperlinkref{2602340166667619136}{Doctest showcase}
\item \hyperlinkref{14134642751930898430}{at-raw blocks}
\item \hyperlinkref{12734752257024730942}{Cross-references}
\item \hyperlinkref{17095889498970318110}{Contents}
\end{itemize}
\end{itemize}


\section{Basic Markdown}



\label{13015211943483912510}{}


Documenter can render all the \href{https://docs.julialang.org/en/v1/stdlib/Markdown/}{Markdown syntax supported by the Julia Markdown parser}. You can use all the usual markdown syntax, such as \textbf{bold text} and \emph{italic text} and \texttt{print({\textquotedbl}inline code{\textquotedbl})}.



\section{Code blocks}



\label{14660894509460337656}{}


Code blocks are rendered as follows:




\begin{minted}{text}
This is an non-highlighted code block.
... Rendered in monospace.
\end{minted}



When the language is specified for the block, e.g. by starting the block with \texttt{```julia}, the contents gets highlighted appropriately (for the language that are supported by the highlighter).




\begin{minted}{julia}
function foo(x::Integer)
    @show x + 1
end
\end{minted}



\section{Mathematics}



\label{9406902154236239312}{}


For mathematics, both inline and display equations are available. Inline equations should be written as LaTeX between two backticks, e.g. \texttt{``A x{\textasciicircum}2 + B x + C = 0``}. It will render as \(A x^2 + B x + C = 0\).



The LaTeX for display equations must be wrapped in a \texttt{```math} code block and will render like



\begin{equation*}
\begin{split}x_{1,2} = \frac{-B \pm \sqrt{B^2 - 4 A C}}{2A}\end{split}\end{equation*}


By default, the HTML output renders equations with \href{https://katex.org/}{KaTeX}, but \href{https://www.mathjax.org/}{MathJax} can optionally be used as well.



\begin{tcolorbox}[toptitle=-1mm,bottomtitle=1mm,colback=admonition-warning!50!white,colframe=admonition-warning,title=\textbf{Warning}]
Similar to LaTeX, using \texttt{\$} and \texttt{\$\$} to escape inline and display equations also works. However, doing so is deprecated and this functionality may be removed in a future release.

\end{tcolorbox}


\section{Images}



\label{5224713065671704740}{}


Include images using basic Markdown syntax:



\begin{figure}
\centering
\includegraphics[max width=\linewidth]{logo.png}
\caption{Enter a descriptive caption for the image}
\end{figure}




The path should be relative to the directory of the current file. Alternatively, use \texttt{./} to begin a path relative to the \texttt{src} of the documents, e.g., \texttt{./logo.png}.



\section{Admonitions}



\label{1916796393499649241}{}


Admonitions are colorful boxes used to highlight parts of the documentation.



Each admonition begins with three \texttt{!!!}, and then the content is indented underneath with four spaces:




\begin{minted}{text}
!!! note "An optional title"
    Here is something you should pay attention to.
\end{minted}



Documenter supports a range of admonition types for different circumstances.



\subparagraph{Note admonition}\indent



\label{3427999462323267970}{}


\begin{tcolorbox}[toptitle=-1mm,bottomtitle=1mm,colback=admonition-note!50!white,colframe=admonition-note,title=\textbf{{\textquotesingle}note{\textquotesingle} admonition}]
Admonitions look like this. This is a \texttt{!!! note}-type admonition.

Note that admonitions themselves can contain other block-level elements too, such as code blocks. E.g.


\begin{minted}{julia}
f(x) = x^2
\end{minted}

However, you \textbf{can not} have at-blocks, docstrings, doctests etc. in an admonition.

Headings are OK though:

\chapter{Heading 1}

\section{Heading 2}

\subsection{Heading 3}

\subsubsection{Heading 4}

\paragraph{Heading 5}\indent

\subparagraph{Heading 6}\indent

\end{tcolorbox}


\subparagraph{Info admonition}\indent



\label{3520455154602929669}{}


\begin{tcolorbox}[toptitle=-1mm,bottomtitle=1mm,colback=admonition-info!50!white,colframe=admonition-info,title=\textbf{{\textquotesingle}info{\textquotesingle} admonition}]
This is a \texttt{!!! info}-type admonition. This is the same as a \texttt{!!! note}-type.

\end{tcolorbox}


\subparagraph{Tip admonition}\indent



\label{8478099471110135607}{}


\begin{tcolorbox}[toptitle=-1mm,bottomtitle=1mm,colback=admonition-tip!50!white,colframe=admonition-tip,title=\textbf{{\textquotesingle}tip{\textquotesingle} admonition}]
This is a \texttt{!!! tip}-type admonition.

\end{tcolorbox}


\subparagraph{Warning admonition}\indent



\label{16113256282093947986}{}


\begin{tcolorbox}[toptitle=-1mm,bottomtitle=1mm,colback=admonition-warning!50!white,colframe=admonition-warning,title=\textbf{{\textquotesingle}warning{\textquotesingle} admonition}]
This is a \texttt{!!! warning}-type admonition.

\end{tcolorbox}


\subparagraph{Danger admonition}\indent



\label{3138459326559650208}{}


\begin{tcolorbox}[toptitle=-1mm,bottomtitle=1mm,colback=admonition-danger!50!white,colframe=admonition-danger,title=\textbf{{\textquotesingle}danger{\textquotesingle} admonition}]
This is a \texttt{!!! danger}-type admonition.

\end{tcolorbox}


\subparagraph{Compat admonition}\indent



\label{8437579397049066350}{}


\begin{tcolorbox}[toptitle=-1mm,bottomtitle=1mm,colback=admonition-compat!50!white,colframe=admonition-compat,title=\textbf{{\textquotesingle}compat{\textquotesingle} admonition}]
This is a \texttt{!!! compat}-type admonition.

\end{tcolorbox}


\subparagraph{Unknown admonition class}\indent



\label{5611698449794175315}{}


\begin{tcolorbox}[toptitle=-1mm,bottomtitle=1mm,colback=admonition-default!50!white,colframe=admonition-default,title=\textbf{Unknown admonition class}]
Admonition with an unknown admonition class. This is a \texttt{code example}.

\end{tcolorbox}


\section{Lists}



\label{932242378642137797}{}


Tight lists look as follows



\begin{itemize}
\item Lorem ipsum dolor sit amet, consectetur adipiscing elit.


\item Nulla quis venenatis justo.


\item In non \emph{sodales} eros.

\end{itemize}


If the lists contain paragraphs or other block level elements, they look like this:



\begin{itemize}
\item Morbi et varius nisl, eu semper orci.

Donec vel nibh sapien. Maecenas ultricies mauris sapien. Nunc et sem ac justo ultricies dignissim ac vitae sem.


\item Nulla molestie aliquet metus, a dapibus ligula.

Morbi pellentesque sodales sollicitudin. Fusce semper placerat suscipit. Aliquam semper tempus ex, non efficitur erat posuere in. Fusce at orci eu ex sagittis commodo.

\begin{quote}
Fusce tempus scelerisque egestas. Pellentesque varius nulla a varius fringilla.

\end{quote}
Fusce nec urna eu orci porta blandit.

\end{itemize}


Numbered lists are also supported



\begin{itemize}
\item[1. ] Lorem ipsum dolor sit amet, consectetur adipiscing elit.


\item[2. ] Nulla quis venenatis justo.


\item[3. ] In non \emph{sodales} eros.

\end{itemize}


As are nested lists



\begin{itemize}
\item Morbi et varius nisl, eu semper orci.

Donec vel nibh sapien. Maecenas ultricies mauris sapien. Nunc et sem ac justo ultricies dignissim ac vitae sem.

\begin{itemize}
\item Lorem ipsum dolor sit amet, consectetur adipiscing elit.


\item Nulla quis venenatis justo.


\item In non \emph{sodales} eros.

\end{itemize}

\item Nulla molestie aliquet metus, a dapibus ligula.

\begin{itemize}
\item[1. ] Lorem ipsum dolor sit amet, consectetur adipiscing elit.


\item[2. ] Nulla quis venenatis justo.


\item[3. ] In non \emph{sodales} eros.

\end{itemize}
Fusce nec urna eu orci porta blandit.

\end{itemize}


Lists can also be included in other blocks that can contain block level items



\begin{tcolorbox}[toptitle=-1mm,bottomtitle=1mm,colback=admonition-note!50!white,colframe=admonition-note,title=\textbf{Bulleted lists in admonitions}]
\begin{itemize}
\item Lorem ipsum dolor sit amet, consectetur adipiscing elit.


\item Nulla quis venenatis justo.


\item In non \emph{sodales} eros.

\end{itemize}
\end{tcolorbox}


\begin{tcolorbox}[toptitle=-1mm,bottomtitle=1mm,colback=admonition-note!50!white,colframe=admonition-note,title=\textbf{Large lists in admonitions}]
\begin{itemize}
\item Morbi et varius nisl, eu semper orci.

Donec vel nibh sapien. Maecenas ultricies mauris sapien. Nunc et sem ac justo ultricies dignissim ac vitae sem.

\begin{itemize}
\item Lorem ipsum dolor sit amet, consectetur adipiscing elit.


\item Nulla quis venenatis justo.


\item In non \emph{sodales} eros.

\end{itemize}

\item Nulla molestie aliquet metus, a dapibus ligula.

\begin{itemize}
\item[1. ] Lorem ipsum dolor sit amet, consectetur adipiscing elit.


\item[2. ] Nulla quis venenatis justo.


\item[3. ] In non \emph{sodales} eros.

\end{itemize}
Fusce nec urna eu orci porta blandit.

\end{itemize}
\end{tcolorbox}


\begin{quote}
\begin{itemize}
\item Morbi et varius nisl, eu semper orci.

Donec vel nibh sapien. Maecenas ultricies mauris sapien. Nunc et sem ac justo ultricies dignissim ac vitae sem.

\begin{itemize}
\item Lorem ipsum dolor sit amet, consectetur adipiscing elit.


\item Nulla quis venenatis justo.


\item In non \emph{sodales} eros.

\end{itemize}
\end{itemize}
\end{quote}


\section{Tables}



\label{6147076236960597988}{}



\begin{table}[h]
\centering
\begin{tabulary}{\linewidth}{R R R}
\toprule
object & implemented & value \\
\toprule
\texttt{A} & ✓ & 10.00 \\
\texttt{BB} & ✓ & 1000000.00 \\
\bottomrule
\end{tabulary}

\end{table}



With explicit alignment.




\begin{table}[h]
\centering
\begin{tabulary}{\linewidth}{L C R}
\toprule
object & implemented & value \\
\toprule
\texttt{A} & ✓ & 10.00 \\
\texttt{BB} & ✓ & 1000000.00 \\
\texttt{CC} & v & 1000000.00 \\
\bottomrule
\end{tabulary}

\end{table}



Huges tables get scaled down:




\begin{table}[h]
\centering
\begin{tabulary}{\linewidth}{C C C C C C C C C C C C C C C C C C C C C C C C C C C C C C C C C C C C C C C C C C C C C C C C C C}
\toprule
1-1 & 1-2 & 1-3 & 1-4 & 1-5 & 1-6 & 1-7 & 1-8 & 1-9 & 1-10 & 1-11 & 1-12 & 1-13 & 1-14 & 1-15 & 1-16 & 1-17 & 1-18 & 1-19 & 1-20 & 1-21 & 1-22 & 1-23 & 1-24 & 1-25 & 1-26 & 1-27 & 1-28 & 1-29 & 1-30 & 1-31 & 1-32 & 1-33 & 1-34 & 1-35 & 1-36 & 1-37 & 1-38 & 1-39 & 1-40 & 1-41 & 1-42 & 1-43 & 1-44 & 1-45 & 1-46 & 1-47 & 1-48 & 1-49 & 1-50 \\
\toprule
2-1 & 2-2 & 2-3 & 2-4 & 2-5 & 2-6 & 2-7 & 2-8 & 2-9 & 2-10 & 2-11 & 2-12 & 2-13 & 2-14 & 2-15 & 2-16 & 2-17 & 2-18 & 2-19 & 2-20 & 2-21 & 2-22 & 2-23 & 2-24 & 2-25 & 2-26 & 2-27 & 2-28 & 2-29 & 2-30 & 2-31 & 2-32 & 2-33 & 2-34 & 2-35 & 2-36 & 2-37 & 2-38 & 2-39 & 2-40 & 2-41 & 2-42 & 2-43 & 2-44 & 2-45 & 2-46 & 2-47 & 2-48 & 2-49 & 2-50 \\
3-1 & 3-2 & 3-3 & 3-4 & 3-5 & 3-6 & 3-7 & 3-8 & 3-9 & 3-10 & 3-11 & 3-12 & 3-13 & 3-14 & 3-15 & 3-16 & 3-17 & 3-18 & 3-19 & 3-20 & 3-21 & 3-22 & 3-23 & 3-24 & 3-25 & 3-26 & 3-27 & 3-28 & 3-29 & 3-30 & 3-31 & 3-32 & 3-33 & 3-34 & 3-35 & 3-36 & 3-37 & 3-38 & 3-39 & 3-40 & 3-41 & 3-42 & 3-43 & 3-44 & 3-45 & 3-46 & 3-47 & 3-48 & 3-49 & 3-50 \\
4-1 & 4-2 & 4-3 & 4-4 & 4-5 & 4-6 & 4-7 & 4-8 & 4-9 & 4-10 & 4-11 & 4-12 & 4-13 & 4-14 & 4-15 & 4-16 & 4-17 & 4-18 & 4-19 & 4-20 & 4-21 & 4-22 & 4-23 & 4-24 & 4-25 & 4-26 & 4-27 & 4-28 & 4-29 & 4-30 & 4-31 & 4-32 & 4-33 & 4-34 & 4-35 & 4-36 & 4-37 & 4-38 & 4-39 & 4-40 & 4-41 & 4-42 & 4-43 & 4-44 & 4-45 & 4-46 & 4-47 & 4-48 & 4-49 & 4-50 \\
5-1 & 5-2 & 5-3 & 5-4 & 5-5 & 5-6 & 5-7 & 5-8 & 5-9 & 5-10 & 5-11 & 5-12 & 5-13 & 5-14 & 5-15 & 5-16 & 5-17 & 5-18 & 5-19 & 5-20 & 5-21 & 5-22 & 5-23 & 5-24 & 5-25 & 5-26 & 5-27 & 5-28 & 5-29 & 5-30 & 5-31 & 5-32 & 5-33 & 5-34 & 5-35 & 5-36 & 5-37 & 5-38 & 5-39 & 5-40 & 5-41 & 5-42 & 5-43 & 5-44 & 5-45 & 5-46 & 5-47 & 5-48 & 5-49 & 5-50 \\
6-1 & 6-2 & 6-3 & 6-4 & 6-5 & 6-6 & 6-7 & 6-8 & 6-9 & 6-10 & 6-11 & 6-12 & 6-13 & 6-14 & 6-15 & 6-16 & 6-17 & 6-18 & 6-19 & 6-20 & 6-21 & 6-22 & 6-23 & 6-24 & 6-25 & 6-26 & 6-27 & 6-28 & 6-29 & 6-30 & 6-31 & 6-32 & 6-33 & 6-34 & 6-35 & 6-36 & 6-37 & 6-38 & 6-39 & 6-40 & 6-41 & 6-42 & 6-43 & 6-44 & 6-45 & 6-46 & 6-47 & 6-48 & 6-49 & 6-50 \\
7-1 & 7-2 & 7-3 & 7-4 & 7-5 & 7-6 & 7-7 & 7-8 & 7-9 & 7-10 & 7-11 & 7-12 & 7-13 & 7-14 & 7-15 & 7-16 & 7-17 & 7-18 & 7-19 & 7-20 & 7-21 & 7-22 & 7-23 & 7-24 & 7-25 & 7-26 & 7-27 & 7-28 & 7-29 & 7-30 & 7-31 & 7-32 & 7-33 & 7-34 & 7-35 & 7-36 & 7-37 & 7-38 & 7-39 & 7-40 & 7-41 & 7-42 & 7-43 & 7-44 & 7-45 & 7-46 & 7-47 & 7-48 & 7-49 & 7-50 \\
8-1 & 8-2 & 8-3 & 8-4 & 8-5 & 8-6 & 8-7 & 8-8 & 8-9 & 8-10 & 8-11 & 8-12 & 8-13 & 8-14 & 8-15 & 8-16 & 8-17 & 8-18 & 8-19 & 8-20 & 8-21 & 8-22 & 8-23 & 8-24 & 8-25 & 8-26 & 8-27 & 8-28 & 8-29 & 8-30 & 8-31 & 8-32 & 8-33 & 8-34 & 8-35 & 8-36 & 8-37 & 8-38 & 8-39 & 8-40 & 8-41 & 8-42 & 8-43 & 8-44 & 8-45 & 8-46 & 8-47 & 8-48 & 8-49 & 8-50 \\
9-1 & 9-2 & 9-3 & 9-4 & 9-5 & 9-6 & 9-7 & 9-8 & 9-9 & 9-10 & 9-11 & 9-12 & 9-13 & 9-14 & 9-15 & 9-16 & 9-17 & 9-18 & 9-19 & 9-20 & 9-21 & 9-22 & 9-23 & 9-24 & 9-25 & 9-26 & 9-27 & 9-28 & 9-29 & 9-30 & 9-31 & 9-32 & 9-33 & 9-34 & 9-35 & 9-36 & 9-37 & 9-38 & 9-39 & 9-40 & 9-41 & 9-42 & 9-43 & 9-44 & 9-45 & 9-46 & 9-47 & 9-48 & 9-49 & 9-50 \\
10-1 & 10-2 & 10-3 & 10-4 & 10-5 & 10-6 & 10-7 & 10-8 & 10-9 & 10-10 & 10-11 & 10-12 & 10-13 & 10-14 & 10-15 & 10-16 & 10-17 & 10-18 & 10-19 & 10-20 & 10-21 & 10-22 & 10-23 & 10-24 & 10-25 & 10-26 & 10-27 & 10-28 & 10-29 & 10-30 & 10-31 & 10-32 & 10-33 & 10-34 & 10-35 & 10-36 & 10-37 & 10-38 & 10-39 & 10-40 & 10-41 & 10-42 & 10-43 & 10-44 & 10-45 & 10-46 & 10-47 & 10-48 & 10-49 & 10-50 \\
\bottomrule
\end{tabulary}

\end{table}



Including ones generated by at-example blocks with \texttt{show} methods:




\begin{minted}{julia}
using Markdown
nrows, ncols = 5, 5
table = Markdown.Table(
    [[[string(i, "-", j)] for j = 1:ncols] for i = 1:nrows],
    [:c for _ in 1:ncols],
)
Markdown.MD([table])
\end{minted}

\begin{table}[h]
\centering
\adjustbox{max width=\linewidth}{
\begin{tabular}
{c | c | c | c | c}
1-1 & 1-2 & 1-3 & 1-4 & 1-5 \\
\hline
2-1 & 2-2 & 2-3 & 2-4 & 2-5 \\
3-1 & 3-2 & 3-3 & 3-4 & 3-5 \\
4-1 & 4-2 & 4-3 & 4-4 & 4-5 \\
5-1 & 5-2 & 5-3 & 5-4 & 5-5 \\
\end{tabular}
}\end{table}



\begin{minted}{julia}
using Markdown
nrows, ncols = 10, 50
table = Markdown.Table(
    [[[string(i, "-", j)] for j = 1:ncols] for i = 1:nrows],
    [:c for _ in 1:ncols],
)
Markdown.MD([table])
\end{minted}

\begin{table}[h]
\centering
\adjustbox{max width=\linewidth}{
\begin{tabular}
{c | c | c | c | c | c | c | c | c | c | c | c | c | c | c | c | c | c | c | c | c | c | c | c | c | c | c | c | c | c | c | c | c | c | c | c | c | c | c | c | c | c | c | c | c | c | c | c | c | c}
1-1 & 1-2 & 1-3 & 1-4 & 1-5 & 1-6 & 1-7 & 1-8 & 1-9 & 1-10 & 1-11 & 1-12 & 1-13 & 1-14 & 1-15 & 1-16 & 1-17 & 1-18 & 1-19 & 1-20 & 1-21 & 1-22 & 1-23 & 1-24 & 1-25 & 1-26 & 1-27 & 1-28 & 1-29 & 1-30 & 1-31 & 1-32 & 1-33 & 1-34 & 1-35 & 1-36 & 1-37 & 1-38 & 1-39 & 1-40 & 1-41 & 1-42 & 1-43 & 1-44 & 1-45 & 1-46 & 1-47 & 1-48 & 1-49 & 1-50 \\
\hline
2-1 & 2-2 & 2-3 & 2-4 & 2-5 & 2-6 & 2-7 & 2-8 & 2-9 & 2-10 & 2-11 & 2-12 & 2-13 & 2-14 & 2-15 & 2-16 & 2-17 & 2-18 & 2-19 & 2-20 & 2-21 & 2-22 & 2-23 & 2-24 & 2-25 & 2-26 & 2-27 & 2-28 & 2-29 & 2-30 & 2-31 & 2-32 & 2-33 & 2-34 & 2-35 & 2-36 & 2-37 & 2-38 & 2-39 & 2-40 & 2-41 & 2-42 & 2-43 & 2-44 & 2-45 & 2-46 & 2-47 & 2-48 & 2-49 & 2-50 \\
3-1 & 3-2 & 3-3 & 3-4 & 3-5 & 3-6 & 3-7 & 3-8 & 3-9 & 3-10 & 3-11 & 3-12 & 3-13 & 3-14 & 3-15 & 3-16 & 3-17 & 3-18 & 3-19 & 3-20 & 3-21 & 3-22 & 3-23 & 3-24 & 3-25 & 3-26 & 3-27 & 3-28 & 3-29 & 3-30 & 3-31 & 3-32 & 3-33 & 3-34 & 3-35 & 3-36 & 3-37 & 3-38 & 3-39 & 3-40 & 3-41 & 3-42 & 3-43 & 3-44 & 3-45 & 3-46 & 3-47 & 3-48 & 3-49 & 3-50 \\
4-1 & 4-2 & 4-3 & 4-4 & 4-5 & 4-6 & 4-7 & 4-8 & 4-9 & 4-10 & 4-11 & 4-12 & 4-13 & 4-14 & 4-15 & 4-16 & 4-17 & 4-18 & 4-19 & 4-20 & 4-21 & 4-22 & 4-23 & 4-24 & 4-25 & 4-26 & 4-27 & 4-28 & 4-29 & 4-30 & 4-31 & 4-32 & 4-33 & 4-34 & 4-35 & 4-36 & 4-37 & 4-38 & 4-39 & 4-40 & 4-41 & 4-42 & 4-43 & 4-44 & 4-45 & 4-46 & 4-47 & 4-48 & 4-49 & 4-50 \\
5-1 & 5-2 & 5-3 & 5-4 & 5-5 & 5-6 & 5-7 & 5-8 & 5-9 & 5-10 & 5-11 & 5-12 & 5-13 & 5-14 & 5-15 & 5-16 & 5-17 & 5-18 & 5-19 & 5-20 & 5-21 & 5-22 & 5-23 & 5-24 & 5-25 & 5-26 & 5-27 & 5-28 & 5-29 & 5-30 & 5-31 & 5-32 & 5-33 & 5-34 & 5-35 & 5-36 & 5-37 & 5-38 & 5-39 & 5-40 & 5-41 & 5-42 & 5-43 & 5-44 & 5-45 & 5-46 & 5-47 & 5-48 & 5-49 & 5-50 \\
6-1 & 6-2 & 6-3 & 6-4 & 6-5 & 6-6 & 6-7 & 6-8 & 6-9 & 6-10 & 6-11 & 6-12 & 6-13 & 6-14 & 6-15 & 6-16 & 6-17 & 6-18 & 6-19 & 6-20 & 6-21 & 6-22 & 6-23 & 6-24 & 6-25 & 6-26 & 6-27 & 6-28 & 6-29 & 6-30 & 6-31 & 6-32 & 6-33 & 6-34 & 6-35 & 6-36 & 6-37 & 6-38 & 6-39 & 6-40 & 6-41 & 6-42 & 6-43 & 6-44 & 6-45 & 6-46 & 6-47 & 6-48 & 6-49 & 6-50 \\
7-1 & 7-2 & 7-3 & 7-4 & 7-5 & 7-6 & 7-7 & 7-8 & 7-9 & 7-10 & 7-11 & 7-12 & 7-13 & 7-14 & 7-15 & 7-16 & 7-17 & 7-18 & 7-19 & 7-20 & 7-21 & 7-22 & 7-23 & 7-24 & 7-25 & 7-26 & 7-27 & 7-28 & 7-29 & 7-30 & 7-31 & 7-32 & 7-33 & 7-34 & 7-35 & 7-36 & 7-37 & 7-38 & 7-39 & 7-40 & 7-41 & 7-42 & 7-43 & 7-44 & 7-45 & 7-46 & 7-47 & 7-48 & 7-49 & 7-50 \\
8-1 & 8-2 & 8-3 & 8-4 & 8-5 & 8-6 & 8-7 & 8-8 & 8-9 & 8-10 & 8-11 & 8-12 & 8-13 & 8-14 & 8-15 & 8-16 & 8-17 & 8-18 & 8-19 & 8-20 & 8-21 & 8-22 & 8-23 & 8-24 & 8-25 & 8-26 & 8-27 & 8-28 & 8-29 & 8-30 & 8-31 & 8-32 & 8-33 & 8-34 & 8-35 & 8-36 & 8-37 & 8-38 & 8-39 & 8-40 & 8-41 & 8-42 & 8-43 & 8-44 & 8-45 & 8-46 & 8-47 & 8-48 & 8-49 & 8-50 \\
9-1 & 9-2 & 9-3 & 9-4 & 9-5 & 9-6 & 9-7 & 9-8 & 9-9 & 9-10 & 9-11 & 9-12 & 9-13 & 9-14 & 9-15 & 9-16 & 9-17 & 9-18 & 9-19 & 9-20 & 9-21 & 9-22 & 9-23 & 9-24 & 9-25 & 9-26 & 9-27 & 9-28 & 9-29 & 9-30 & 9-31 & 9-32 & 9-33 & 9-34 & 9-35 & 9-36 & 9-37 & 9-38 & 9-39 & 9-40 & 9-41 & 9-42 & 9-43 & 9-44 & 9-45 & 9-46 & 9-47 & 9-48 & 9-49 & 9-50 \\
10-1 & 10-2 & 10-3 & 10-4 & 10-5 & 10-6 & 10-7 & 10-8 & 10-9 & 10-10 & 10-11 & 10-12 & 10-13 & 10-14 & 10-15 & 10-16 & 10-17 & 10-18 & 10-19 & 10-20 & 10-21 & 10-22 & 10-23 & 10-24 & 10-25 & 10-26 & 10-27 & 10-28 & 10-29 & 10-30 & 10-31 & 10-32 & 10-33 & 10-34 & 10-35 & 10-36 & 10-37 & 10-38 & 10-39 & 10-40 & 10-41 & 10-42 & 10-43 & 10-44 & 10-45 & 10-46 & 10-47 & 10-48 & 10-49 & 10-50 \\
\end{tabular}
}\end{table}


However, tables with huge cells are not properly handled right now:




\begin{table}[h]
\centering
\begin{tabulary}{\linewidth}{L C R}
\toprule
object & implemented & value \\
\toprule
\texttt{A} & ✓ & 10.00 \\
\texttt{BBBBBBBBBBBBBBBBBBBB} & ✓✓✓✓✓✓✓✓✓✓✓✓✓✓✓✓✓✓✓✓✓✓✓✓✓✓✓✓✓✓✓✓✓✓✓✓✓✓✓✓✓✓✓✓ & 1000000000000000000000000000000000000000000000000000000.00 \\
\bottomrule
\end{tabulary}

\end{table}



\section{Footnotes}



\label{16845496644619434430}{}


Footnote references can be added with the \texttt{[{\textasciicircum}label]} syntax.\footnotemark[1] The footnote definitions get collected at the bottom of the page.



The footnote label can be an arbitrary string and even consist of block-level elements.\footnotemark[2]



\footnotetext[1]{A footnote definition uses the \texttt{[{\textasciicircum}label]: ...} syntax in a block scope.

}


\footnotetext[2]{\begin{quote}
Any sufficiently advanced technology is indistinguishable from magic.

\end{quote}
Arthur C. Clarke, \emph{Profiles of the Future} (1961): Clarke{\textquotesingle}s Third Law.

}


Footnote with more block contents.\footnotemark[3]



\footnotetext[3]{And a link to another footnote\footnotemark[1].

\begin{tcolorbox}[toptitle=-1mm,bottomtitle=1mm,colback=admonition-note!50!white,colframe=admonition-note,title=\textbf{Admonition..}]
.. in a footnote.

\end{tcolorbox}
}


\section{Headings}



\label{5680510222270408269}{}


Finally, headings render as follows



\subsection{Heading level 3}



\label{13085847155693612023}{}


\subsubsection{Heading level 4}



\label{12782044777122551513}{}


\paragraph{Heading level 5}\indent



\label{9543099701887315540}{}


\subparagraph{Heading level 6}\indent



\label{10897944457118690625}{}


To see an example of a level 1 heading see the page title and for level 2 heading, see the one just under this paragraph.



\begin{tcolorbox}[toptitle=-1mm,bottomtitle=1mm,colback=admonition-note!50!white,colframe=admonition-note,title=\textbf{Headings in sidebars}]
Level 1 and 2 heading show up in the sidebar, for the current page.

\end{tcolorbox}


\section{Doctesting example}



\label{10597583416623881668}{}


Often you want to write code example such as this:




\begin{minted}{jlcon}
julia> f(x) = x^2
f (generic function with 1 method)

julia> f(3)
9
\end{minted}



If you write them as a \texttt{```jldoctest} code block, Documenter can make sure that the doctest has not become outdated. See \href{@ref}{Doctests} for more information.



Script-style doctests are supported too:




\begin{minted}{julia}
2 + 2
# output
4
\end{minted}



\section{Running interactive code}



\label{3545462789522369084}{}


\hyperlinkref{12317204655950454097}{\texttt{@example} block} run a code snippet and insert the output into the document. E.g. the following Markdown




\begin{minted}{text}
```@example
typeof(π)
```
\end{minted}



becomes the following code-output block pair




\begin{minted}{julia}
typeof(π)
\end{minted}


\begin{minted}[xleftmargin=-\fboxsep,xrightmargin=-\fboxsep,bgcolor=white,frame=single]{text}
Irrational{:π}
\end{minted}



If the last element can be rendered as \texttt{text/latex}:




\begin{minted}{julia}
struct LaTeXRender
    s :: String
end
function Base.show(io, ::MIME"text/latex", r::LaTeXRender)
    write(io, """
    Rendering \\texttt{LaTeXRender}: $(r.s)
    """)
end
\end{minted}




\begin{minted}{julia}
LaTeXRender("render this")
\end{minted}

Rendering \texttt{LaTeXRender}: render this


If the last value in an \texttt{@example} block is a \texttt{nothing}, the standard output from the blocks{\textquotesingle} evaluation gets displayed instead




\begin{minted}{julia}
println("Hello World")
\end{minted}


\begin{minted}[xleftmargin=-\fboxsep,xrightmargin=-\fboxsep,bgcolor=white,frame=single]{text}
Hello World
\end{minted}



However, do note that if the block prints to standard output, but also has a final non-\texttt{nothing} value, the standard output just gets discarded:




\begin{minted}{julia}
println("Hello World")
42
\end{minted}


\begin{minted}[xleftmargin=-\fboxsep,xrightmargin=-\fboxsep,bgcolor=white,frame=single]{text}
42
\end{minted}



\subsection{at-REPL blocks}



\label{9794080512499797594}{}



\begin{minted}{jlcon}
julia> (1 // 2)^2
1//4
\end{minted}




\begin{minted}{jlcon}
julia> x = 1
1

julia> x // 2
1//2

julia> ans^2
1//4
\end{minted}



\section{Setup blocks}



\label{5969831508618549664}{}


\texttt{SetupNode} with \texttt{x = 5}...





...makes the following evaluate to...




\begin{minted}{julia}
x + 1
\end{minted}


\begin{minted}[xleftmargin=-\fboxsep,xrightmargin=-\fboxsep,bgcolor=white,frame=single]{text}
6
\end{minted}



\subsection{Color output}



\label{17347250707405175741}{}


Output from \hyperlinkref{16839343392674353047}{\texttt{@repl} block}s and \hyperlinkref{12317204655950454097}{\texttt{@example} block}s support colored output, transforming ANSI color codes to HTML.



\subsubsection{Colored \texttt{@repl} block output}



\label{681009610151884915}{}


\textbf{Input:}




\begin{minted}{text}
```@repl
printstyled("This should be in bold light cyan.", color=:light_cyan, bold=true)
```
\end{minted}



\textbf{Output:}




\begin{minted}{jlcon}
julia> printstyled("This should be in bold cyan.", color=:cyan, bold=true)
This should be in bold cyan.
\end{minted}



\textbf{Locally disabled color:}




\begin{minted}{text}
```@repl; ansicolor=false
printstyled("This should be in bold light cyan.", color=:light_cyan, bold=true)
```
\end{minted}




\begin{minted}{jlcon}
julia> printstyled("This should be in bold light cyan.", color=:light_cyan, bold=true)
This should be in bold light cyan.
\end{minted}



\subsubsection{Raw ANSI code output}



\label{13267379516757075300}{}


Regardless of the color setting, when you print the ANSI escape codes directly, coloring is enabled.




\begin{minted}{julia}
for color in 0:15
    print("\e[38;5;$color;48;5;$(color)m  ")
    print("\e[49m", lpad(color, 3), " ")
    color % 8 == 7 && println()
end
print("\e[m")
\end{minted}


\begin{minted}[xleftmargin=-\fboxsep,xrightmargin=-\fboxsep,bgcolor=white,frame=single]{text}
    0     1     2     3     4     5     6     7
    8     9    10    11    12    13    14    15

\end{minted}



\section{Doctest showcase}



\label{2602340166667619136}{}


Currently exists just so that there would be doctests to run in manual pages of Documenter{\textquotesingle}s manual. This page does not show up in navigation.




\begin{minted}{jlcon}
julia> 2 + 2
4
\end{minted}



The following doctests needs doctestsetup:




\begin{minted}{jlcon}
julia> Documenter.splitexpr(:(Foo.Bar.baz))
(:(Foo.Bar), :(:baz))
\end{minted}



Let{\textquotesingle}s also try \texttt{@meta} blocks:






\begin{minted}{jlcon}
julia> f(2)
4
\end{minted}





\section{at-raw blocks}



\label{14134642751930898430}{}


Only one of them should end up in the output.






bar



\section{Cross-references}



\label{12734752257024730942}{}


Here we link to the heading \hyperlinkref{12734752257024730942}{Cross-references}. \hyperlinkref{12734752257024730942}{This} also links there.



We can also link to other pages, such as: \hyperlinkref{4901604269105032859}{An index of docstrings}.



Then there are link to docstrings: \href{@ref}{\texttt{Foo.bar}}, \href{@ref Foo.bar}{\texttt{foobar}}, \href{@ref Foo.bar}{foobar}.



And issue references: \href{https://example.org/Repository.jl/blob/1245}{\#1245}



\subsection{\texttt{@example}-block}



\label{12317204655950454097}{}


\subsection{\texttt{@contents}-block}



\label{5495928211144802529}{}


\subsection{\texttt{@repl}-block}



\label{16839343392674353047}{}


\section{Contents}



\label{17095889498970318110}{}

\begin{itemize}
\item \hyperlinkref{4859465668994674752}{Docstrings}
\begin{itemize}
\item \hyperlinkref{4901604269105032859}{An index of docstrings}
\item \hyperlinkref{17848419264812669351}{at-autodocs}
\end{itemize}
\item \hyperlinkref{5852870456629818970}{Style showcase for \LaTeX}
\begin{itemize}
\item \hyperlinkref{11192443102986465861}{Table of contents}
\item \hyperlinkref{13015211943483912510}{Basic Markdown}
\item \hyperlinkref{14660894509460337656}{Code blocks}
\item \hyperlinkref{9406902154236239312}{Mathematics}
\item \hyperlinkref{5224713065671704740}{Images}
\item \hyperlinkref{1916796393499649241}{Admonitions}
\item \hyperlinkref{932242378642137797}{Lists}
\item \hyperlinkref{6147076236960597988}{Tables}
\item \hyperlinkref{16845496644619434430}{Footnotes}
\item \hyperlinkref{5680510222270408269}{Headings}
\item \hyperlinkref{10597583416623881668}{Doctesting example}
\item \hyperlinkref{3545462789522369084}{Running interactive code}
\item \hyperlinkref{5969831508618549664}{Setup blocks}
\item \hyperlinkref{2602340166667619136}{Doctest showcase}
\item \hyperlinkref{14134642751930898430}{at-raw blocks}
\item \hyperlinkref{12734752257024730942}{Cross-references}
\item \hyperlinkref{17095889498970318110}{Contents}
\end{itemize}
\end{itemize}


\chapter{Docstrings}



\label{4859465668994674752}{}


The following example docstrings come from the demo modules in \texttt{make.jl}.


\hypertarget{5669384191785722396}{\texttt{Main.Mod.func}}  -- {Function.}

\begin{adjustwidth}{2em}{0pt}


\begin{minted}{julia}
func(x)
\end{minted}

\hyperlinkref{1885743281855441478}{\texttt{T}}



\href{https://example.org/Repository.jl/blob/make.jl#L30-34}{\texttt{source}}


\end{adjustwidth}
\hypertarget{1885743281855441478}{\texttt{Main.Mod.T}}  -- {Type.}

\begin{adjustwidth}{2em}{0pt}


\begin{minted}{julia}
T
\end{minted}

\hyperlinkref{5669384191785722396}{\texttt{func(x)}}



\href{https://example.org/Repository.jl/blob/make.jl#L37-41}{\texttt{source}}


\end{adjustwidth}

\section{An index of docstrings}



\label{4901604269105032859}{}


The \href{@ref}{\texttt{@index} block} can be used to generate a list of all the docstrings on a page (or even across pages) and will look as follows


\begin{itemize}
\item \hyperlinkref{11302554442225637105}{\texttt{AutoDocs}}
\item \hyperlinkref{15441025252371609530}{\texttt{Main.AutoDocs.K}}
\item \hyperlinkref{12338358314596597808}{\texttt{Main.AutoDocs.T}}
\item \hyperlinkref{1885743281855441478}{\texttt{Main.Mod.T}}
\item \hyperlinkref{14053004641171891989}{\texttt{Main.AutoDocs.f}}
\item \hyperlinkref{5669384191785722396}{\texttt{Main.Mod.func}}
\item \hyperlinkref{15818938160560313688}{\texttt{Main.Mod.long\_equations\_in\_docstrings}}
\item \hyperlinkref{7256264955516068825}{\texttt{Main.AutoDocs.@m}}
\end{itemize}


{\rule{\textwidth}{1pt}}


Missing docstring:



\begin{tcolorbox}[toptitle=-1mm,bottomtitle=1mm,colback=admonition-warning!50!white,colframe=admonition-warning,title=\textbf{Missing docstring.}]
Missing docstring for \texttt{this\_docstring\_does\_not\_exist}. Check Documenter{\textquotesingle}s build log for details.

\end{tcolorbox}

\hypertarget{15818938160560313688}{\texttt{Main.Mod.long\_equations\_in\_docstrings}}  -- {Function.}

\begin{adjustwidth}{2em}{0pt}

Inline: \(\frac{1+2+3+4+5+6}{\sigma^2} + \frac{1+2+3+4+5+6}{\sigma^2} + \frac{1+2+3+4+5+6}{\sigma^2} + \frac{1+2+3+4+5+6}{\sigma^2} + \frac{1+2+3+4+5+6}{\sigma^2} + \frac{1+2+3+4+5+6}{\sigma^2} + \frac{1+2+3+4+5+6}{\sigma^2} + \frac{1+2+3+4+5+6}{\sigma^2}\)

Display equation:

\begin{equation*}
\begin{split}\frac{1+2+3+4+5+6}{\sigma^2} + \frac{1+2+3+4+5+6}{\sigma^2} + \frac{1+2+3+4+5+6}{\sigma^2} + \frac{1+2+3+4+5+6}{\sigma^2} + \frac{1+2+3+4+5+6}{\sigma^2} + \frac{1+2+3+4+5+6}{\sigma^2} + \frac{1+2+3+4+5+6}{\sigma^2} + \frac{1+2+3+4+5+6}{\sigma^2}\end{split}\end{equation*}
\begin{tcolorbox}[toptitle=-1mm,bottomtitle=1mm,colback=admonition-note!50!white,colframe=admonition-note,title=\textbf{Long equations in admonitions}]
Inline: \(\frac{1+2+3+4+5+6}{\sigma^2} + \frac{1+2+3+4+5+6}{\sigma^2} + \frac{1+2+3+4+5+6}{\sigma^2} + \frac{1+2+3+4+5+6}{\sigma^2} + \frac{1+2+3+4+5+6}{\sigma^2} + \frac{1+2+3+4+5+6}{\sigma^2} + \frac{1+2+3+4+5+6}{\sigma^2} + \frac{1+2+3+4+5+6}{\sigma^2}\)

Display equation:

\begin{equation*}
\begin{split}\frac{1+2+3+4+5+6}{\sigma^2} + \frac{1+2+3+4+5+6}{\sigma^2} + \frac{1+2+3+4+5+6}{\sigma^2} + \frac{1+2+3+4+5+6}{\sigma^2} + \frac{1+2+3+4+5+6}{\sigma^2} + \frac{1+2+3+4+5+6}{\sigma^2} + \frac{1+2+3+4+5+6}{\sigma^2} + \frac{1+2+3+4+5+6}{\sigma^2}\end{split}\end{equation*}
\end{tcolorbox}
Long equations in footnotes.\footnotemark[1]

\footnotetext[1]{Inline: \(\frac{1+2+3+4+5+6}{\sigma^2} + \frac{1+2+3+4+5+6}{\sigma^2} + \frac{1+2+3+4+5+6}{\sigma^2} + \frac{1+2+3+4+5+6}{\sigma^2} + \frac{1+2+3+4+5+6}{\sigma^2} + \frac{1+2+3+4+5+6}{\sigma^2} + \frac{1+2+3+4+5+6}{\sigma^2} + \frac{1+2+3+4+5+6}{\sigma^2} + \frac{1+2+3+4+5+6}{\sigma^2} + \frac{1+2+3+4+5+6}{\sigma^2} + \frac{1+2+3+4+5+6}{\sigma^2} + \frac{1+2+3+4+5+6}{\sigma^2}\)

Display equation:

\begin{equation*}
\begin{split}\frac{1+2+3+4+5+6}{\sigma^2} + \frac{1+2+3+4+5+6}{\sigma^2} + \frac{1+2+3+4+5+6}{\sigma^2} + \frac{1+2+3+4+5+6}{\sigma^2} + \frac{1+2+3+4+5+6}{\sigma^2} + \frac{1+2+3+4+5+6}{\sigma^2} + \frac{1+2+3+4+5+6}{\sigma^2} + \frac{1+2+3+4+5+6}{\sigma^2} + \frac{1+2+3+4+5+6}{\sigma^2} + \frac{1+2+3+4+5+6}{\sigma^2} + \frac{1+2+3+4+5+6}{\sigma^2} + \frac{1+2+3+4+5+6}{\sigma^2}\end{split}\end{equation*}
}


\href{https://example.org/Repository.jl/blob/make.jl#L44-74}{\texttt{source}}


\end{adjustwidth}

\section{at-autodocs}



\label{17848419264812669351}{}

\hypertarget{11302554442225637105}{\texttt{AutoDocs}}  -- {Module.}

\begin{adjustwidth}{2em}{0pt}

\texttt{AutoDocs} module.



\href{https://example.org/Repository.jl/blob/make.jl#L78-78}{\texttt{source}}


\end{adjustwidth}
\hypertarget{15441025252371609530}{\texttt{Main.AutoDocs.K}}  -- {Constant.}

\begin{adjustwidth}{2em}{0pt}

Constant \texttt{K}.



\href{https://example.org/Repository.jl/blob/make.jl#L91-91}{\texttt{source}}


\end{adjustwidth}
\hypertarget{12338358314596597808}{\texttt{Main.AutoDocs.T}}  -- {Type.}

\begin{adjustwidth}{2em}{0pt}

Type \texttt{T}.



\href{https://example.org/Repository.jl/blob/make.jl#L94-94}{\texttt{source}}


\end{adjustwidth}
\hypertarget{14053004641171891989}{\texttt{Main.AutoDocs.f}}  -- {Method.}

\begin{adjustwidth}{2em}{0pt}

Function \texttt{f}.



\href{https://example.org/Repository.jl/blob/make.jl#L88-88}{\texttt{source}}


\end{adjustwidth}
\hypertarget{7256264955516068825}{\texttt{Main.AutoDocs.@m}}  -- {Macro.}

\begin{adjustwidth}{2em}{0pt}

Macro \texttt{@m}.



\href{https://example.org/Repository.jl/blob/make.jl#L97-97}{\texttt{source}}


\end{adjustwidth}

\end{document}
